\chapter{Conclusion}
This master thesis was motivated by the increasing concern over the rising number of car accidents caused by driver distraction and the need for advanced action recognition systems to promote road safety. The main purpose of this thesis was to enhance vehicle interior action recognition using self-supervised learning with 3D human skeleton representations. By utilizing unannotated data, self-supervised learning offers a promising approach to learning meaningful representations for action recognition without the need for labor-intensive annotations.

To achieve our research objectives, we investigated two different approaches for SSL: DINO with STTFormer and PSTL with ST-GCN. The second approach showed promising results in extracting meaningful features from 3D human skeleton representations and demonstrated the potential of self-supervised learning for action recognition in the car interior. Throughout our investigation, we carefully selected evaluation protocols and conducted comprehensive comparisons to answer our research question about recognizing human action based on skeleton data using self-supervised learning. Our experimental results showed that PSTL with ST-GCN achieved competitive performance in action recognition on the NTU RGB+D and the approach DINO with STTFormer has its limitations and needs further research. 

The application of DINO with STTFormer to the task of skeleton-based action recognition encountered certain limitations that hindered its performance. One possible reason for its suboptimal results lies in the fact that the DINO framework has primarily showcased strong performance on image data, which inherently contains a vast amount of information due to the sheer number of pixels. In contrast, skeleton data represents a compressed form of human action, encapsulating movement patterns with fewer dimensions. This limitation also extends to STTFormer, a variant of Transformers that may require a larger volume of data to achieve optimal performance. Additionally, the DINO framework was initially designed for static images, while our task involves sequential data in the form of skeleton sequences. These factors collectively suggest that DINO's efficacy may be hampered when applied directly to skeleton-based action recognition. To further validate the potential of DINO, future research could explore alternative encoders, such as ST-GCN, within the DINO framework. This investigation would provide a more comprehensive understanding of DINO's adaptability and effectiveness for our specific task.

During our exploration of the Drive\&Act dataset, we encountered several challenges, including errors in annotating the dataset, issues with missing or incorrect joint coordinates, and the presence of highly imbalanced classes. To overcome these challenges and achieve better results, future work is needed to develop improved techniques for dealing with the skeleton modality in the Drive\&Act dataset. Additionally, there is an urgent need to create more high-quality datasets dedicated to human action recognition inside the car, which will facilitate the development and evaluation of more accurate and robust action recognition systems.

Furthermore, an area of future research lies in improving the interpretability of self-supervised learning models. As these models learn representations in a self-supervised manner, understanding the learned features and how they correspond to specific human actions and behaviors remains a valuable avenue for exploration.

Throughout the course of this master's thesis, I have delved into the realm of contrastive self-supervised learning for action recognition using 3D skeleton data, a topic that continues to be at the forefront of research and exploration. This endeavor has illuminated the promising potential of applying such techniques to diverse domains, including the recognition of human actions. As I delved into the intricacies of analyzing and evaluating skeleton data, I gained valuable insights into the nuances of this unique data type, which diverges from the more commonly encountered image data. Moreover, this thesis provided an invaluable opportunity to learn how to construct intricate models and architectures, harnessing the power of parallel GPUs to accelerate training processes on server setups.

The journey through this research venture has been a master class in evaluation methodologies. The mastery of diverse evaluation techniques, ranging from k-NN and Linear evaluation to finetuning and cross-evaluation, was particularly enlightening. Each of these approaches granted me a multifaceted perspective on the performance of the proposed methods, enabling comprehensive insights into their strengths and limitations.Another pivotal aspect of this master's thesis was the experience of dealing with various datasets, each accompanied by its distinct set of challenges. This hands-on experience underscored the importance of data preprocessing and augmentation, as well as the need to carefully design experiments to achieve meaningful and reliable results. 


In conclusion, this master thesis has provided valuable insights into the potential of self-supervised learning for human action recognition in the car interior. By leveraging 3D human skeleton representations and employing contrastive learning techniques, our study contributes to the advancement of action recognition systems for improving road safety and driver assistance in both manual and autonomous driving scenarios. The future of self-supervised learning in this domain looks promising, and continued research in this field is vital to unlocking the full potential of action recognition systems for enhanced vehicle safety.
